\documentclass[11pt]{article}
\usepackage[T1]{fontenc}
\usepackage{lmodern}
\usepackage[margin=0.9in]{geometry}
\usepackage{amsmath,amssymb,amsthm,mathtools,array,booktabs}
\usepackage[hidelinks]{hyperref}
\newtheorem{theorem}{Theorem}[section]
\newtheorem{lemma}[theorem]{Lemma}
\newtheorem{corollary}[theorem]{Corollary}
\setlength{\parskip}{0.25em}
\setlength{\parindent}{1.2em}
\allowdisplaybreaks
\title{Improved Guarantees for Online Consensus Clustering}
\author{Research notes accompanying the current draft}
\date{November 16, 2026}
\begin{document}
\maketitle

\section{Scope and main improvements}

These notes give improvements to the current \texttt{main.tex}, with
full proofs of the new inequalities and reductions. They use its
input-label access model and its verified implicit Pivot maintenance
lemmas. All input streams are oblivious. These notes are provided
separately; no changes to existing manuscript files were made in
preparing them.

For partitions $C,P$ on a vertex set $V$, let $d(C,P)$ count unordered
pairs whose together/separate decisions differ. Write
\[
 \operatorname{cost}_t(C)=\frac1m\sum_{i=1}^m d(C,P_i|_{V_t}),
 \qquad \operatorname{OPT}_t=\min_C\operatorname{cost}_t(C).
\]
The offline notation omits $t$. Restricting a partition to fewer vertices
cannot increase its cost. The pair-distance triangle inequality implies
$d(C,F)\le\operatorname{cost}_t(C)+\operatorname{cost}_t(F)$.
In the online statements, $\eta\ge1$ is the factor of the fixed
deterministic Online Agreement solver with $O_\eta(\log n)$ per-vertex
recourse, as in the draft and \cite{online}. It is not an assertion
about the best possible factor of all online solvers.

\begin{center}
\small
\renewcommand{\arraystretch}{1.35}
\begin{tabular}{@{}>{\raggedright\arraybackslash}p{0.24\textwidth}>{\raggedright\arraybackslash}p{0.30\textwidth}>{\raggedright\arraybackslash}p{0.37\textwidth}@{}}
\toprule
Quantity & Current draft & Improvement proved here\\
\midrule
Online competitive ratio
 & $2-\dfrac{2}{3\eta+5}+\varepsilon$
 & $2-\dfrac1{\eta+1}+\varepsilon$\\
Deterministic online recourse
 & $O_\eta(\varepsilon^{-2}\log(2n/\zeta)\log(2mn/\zeta))$
 & $O_\eta(\varepsilon^{-2}\log(2n/\zeta)\log(2n/(\varepsilon\zeta)))$\\
Offline and implicit approximation
 & $1.6944+\varepsilon$
 & $1.64+\varepsilon$ with the original seven samples\\
Offline and implicit approximation with a larger sample
 & $1.6944+\varepsilon$
 & $11/7+\varepsilon$, with polynomial dependence on $1/\varepsilon$
   (Theorem~\ref{thm:eleven-sevenths})\\
Implicit update time
 & $O(\varepsilon^{-3}\Lambda^2\log^2(N+2))$
 & $O(\varepsilon^{-3}\Lambda^2\log(N+2))$ on a word RAM\\
Membership / same-cluster query
 & $O(\log(n+2))$ worst case
 & $O(1)$ worst case on a word RAM\\
\bottomrule
\end{tabular}
\end{center}
Here $n$ denotes the horizon in online statements and the active size
in query statements; $N$ is the maximum active size, $U$ is the update
bound, and $\Lambda=\log((U+2)/(\varepsilon\zeta))$.
The update bounds are expected amortized. A randomized phase selector
also improves the switching contribution to online recourse from the
number of candidates to its logarithm, in expectation and with high
probability; see Theorem~\ref{thm:improved-random-recourse}.
The table's $\varepsilon^{-3}$ update bound applies to the fixed-sample
variants. The separate $11/7+\varepsilon$ result has
$\varepsilon^{-O(1)}$ dependence, stated explicitly in
Theorem~\ref{thm:eleven-sevenths}.

The same centered charging argument gives $1.6+\varepsilon$ using
fifteen sampled inputs. More substantially, it gives $11/7+\varepsilon$
with \emph{polynomial} dependence on $1/\varepsilon$ in running time
and space, improving the exponential dependence in the previous
version of these notes. For fixed accuracy, the offline time remains
near-linear in $n$ and the implicit dynamic update time remains
polylogarithmic. The value $11/7$ itself is the classical consensus
approximation of Ailon, Charikar, and Newman~\cite{acn}; the result here
obtains its approximation from a sample of
$O(\log(1/\varepsilon))$ input partitions and supports the bucket
implementations. It does not assert a new unrestricted offline
approximation record. The polynomial-accuracy proof has large absolute
constants, including a minimum of 87 sampled inputs.

The seven-sample $1.64+\varepsilon$ result uses exactly the draft's
35 buckets per vertex, improving the previous $1.68+\varepsilon$
bound while reducing its sample and bucket counts. A separate
nine-sample construction gives factor $1.64$ without an accuracy
parameter, eliminating the $\varepsilon^{-3}$ dependence when that
fixed factor suffices. All these offline approximation improvements
combine with the faster implicit implementation below.

The explicit insertion-only transfer in Section~\ref{sec:conditional}
is conditional on the stated correlation-clustering solver hypothesis.
It does not resolve the unproved guarantees for the modified solver
identified in \texttt{review.tex}. This caveat concerns the version
assessed in that review; subsequent edits to the solver files are
outside the scope of these notes.

\section{A stronger online competitive ratio}
\label{sec:improved-online}

Write $c=\eta+1$, where $\eta$ is the approximation factor of the
fixed deterministic Online Agreement algorithm used in the draft.
Its persistent identifiers change at most $q_{\rm A}\lceil\log_2 n\rceil$
times per vertex. The improvement below uses the same black box.
Define
\[
 R_* = 2-\frac{1}{\eta+1}=2-\frac1c.
\]
The improvement over the draft's constant is
\[
 \left(2-\frac{2}{3\eta+5}\right)-R_*
 =\frac{\eta+3}{(\eta+1)(3\eta+5)}>0.
\]
For example, the two expressions are $7/4$ and $3/2$ when $\eta=1$.
This example illustrates the formula and makes no claim that an exact
online correlation-clustering solver is available. As $\eta$ increases,
the improvement below two increases by a factor tending to $3/2$.

\subsection{A mixed majority with a weighted error bound}

Put $r=1/\eta\in(0,1]$. Sample the size $k$ to be five with probability
\[
 a_r=\frac{2339-2048r}{2595}
\]
and fifteen with probability
\[
 b_r=\frac{256+2048r}{2595}=1-a_r.
\]
Both probabilities are positive for $0\le r\le1$.
Independently sample $k$ input indices uniformly with replacement.
Join a pair in the sampled signed graph if a strict majority of these
$k$ input partitions joins it. Retain the sample throughout the
vertex-arrival stream.

\begin{lemma}\label{lem:improved-majority}
For odd $k$, write
\[
 \phi_k(p)=\sum_{j=(k+1)/2}^k\binom{k}{j}p^j(1-p)^{k-j}.
\]
If a fraction $p$ of the inputs disagree with a reference decision on
a pair, the sampled majority disagrees with probability
\[
 h_r(p)=a_r\phi_5(p)+b_r\phi_{15}(p).
\]
For all $0\le p,r\le1$,
\[
 {}[1+r(1-2p)]h_r(p)\le2p^2.
\]
Equivalently, with $r=1/\eta$ and $c=\eta+1$,
\[
 (c-2p)h_{1/\eta}(p)\le2\eta p^2.
\]
\end{lemma}
\begin{proof}
The binomial formula follows from independent sampling with replacement.
We give an exact polynomial certificate for the inequality. Define
$B_{i,d}(x)=\binom di x^i(1-x)^{d-i}$ and let $M$ be the following
matrix, whose rows are indexed by $i=0,\ldots,12$ and columns by
$j=0,1,2$,
\[
 M=\begin{pmatrix}
71933400&71933400&71933400\\
68895750&67215225&89189100\\
66657150&67444125&117690300\\
65955078&76436514&166396734\\
68411924&102548054&254296728\\
77927850&165781315&425523420\\
80242530&183628395&350020980\\
72594970&156728740&262939950\\
63845040&126725060&196199640\\
55750590&101452365&148444380\\
48500550&81470025&114439500\\
41961150&65938950&89916750\\
35966700&53950050&71933400
\end{pmatrix}.
\]
Direct polynomial expansion gives the identity
\[
 2p^2-[1+r(1-2p)]h_r(p)
 =\frac{p^2(2p-1)^2}{35966700}
       \sum_{i=0}^{12}\sum_{j=0}^2 M_{i,j}B_{i,12}(p)B_{j,2}(r).
\]
Every factor on the right is nonnegative on the unit square. The
smallest matrix entry is 35966700, so the Bernstein polynomial after
division by 35966700 is at least one there. This proves the inequality,
including $p=0$ and $p=1/2$.
\end{proof}

\paragraph{A smaller sample when $\eta\ge16/3$.}
For $r=1/\eta\le3/16$, the same weighted inequality holds with the
simpler error function
\[
 \widetilde h_r(p)
   =\frac{3-16r}{5}\phi_5(p)+\frac{2+16r}{5}\phi_7(p).
\]
Thus only five or seven inputs need be sampled in this parameter range.
Both mixture weights are nonnegative, and the exact factorization is
\[
 2p^2-[1+r(1-2p)]\widetilde h_r(p)
 =p^2(2p-1)^2\bigl(Q_0(p)+rQ_1(p)+r^2Q_2(p)\bigr),
\]
where
\begin{align*}
 Q_0(p)&=1+(1-p)(1+3p-2p^2),\\
 Q_1(p)&=p\bigl(1+5(1-p)+16(1-p)^2+4(1-p)^3\bigr),\\
 Q_2(p)&=32p(1-p)^3.
\end{align*}
Each expression is nonnegative on $[0,1]$. This proves the desired
inequality directly and gives the same competitive ratio with smaller
absolute constants for the graph representation.

\subsection{A sharper history comparison}

For a reference partition $F$ of $V_t$, let $p_e$ be the fraction of
input partitions disagreeing with $F$ on pair $e$. Let $H$ be the signed
graph of a persistent majority candidate, and let $B_H^t(F)$ be its
pairs disagreeing with $F$.

\begin{lemma}\label{lem:improved-weighted-history}
For every deterministic Online Agreement candidate $A$ and every
reference partition $F$ of $V_t$,
\[
 \max_{u\le t}\operatorname{cost}_u(A(u))
 \le\operatorname{cost}_t(F)
       +\sum_{e\in B_H^t(F)}(c-2p_e).
\]
More precisely, with $F_u=F|_{V_u}$,
\[
 \operatorname{cost}_u(A(u))-\operatorname{cost}_u(F_u)
 \le\sum_{e\in B_H^t(F)\cap\binom{V_u}{2}}(c-2p_e)
 \qquad(u\le t).
\]
\end{lemma}
\begin{proof}
Fix $u\le t$. For each pair in $V_u$, let $y_e$ indicate that $A(u)$
disagrees with $F_u$, and let $b_e$ indicate that $H$ disagrees with
$F_u$. The consensus-cost difference on this pair is $y_e(1-2p_e)$.
The following inequality holds for both values of $b_e$ and $y_e$,
\[
 y_e(1-2p_e)\le \mathbf1[y_e\ne b_e]+b_e(1-2p_e).
\]
Indeed, when $b_e=0$ it is $y_e(1-2p_e)\le y_e$. When $b_e=1$,
the case $y_e=1$ is equality and the case $y_e=0$ follows from
$2-2p_e\ge0$. Summing and using the solver guarantee against $F_u$
gives
\begin{align*}
 \operatorname{cost}_u(A(u))-\operatorname{cost}_u(F_u)
 &\le \operatorname{err}_{H|_{V_u}}(A(u))
          +\sum_e b_e(1-2p_e)\\
 &\le \eta\sum_e b_e+\sum_e b_e(1-2p_e)
 =\sum_{e\in B_H^t(F)\cap\binom{V_u}{2}}(c-2p_e).
\end{align*}
Here $F_u$ is a feasible correlation clustering with error $\sum_e b_e$.
The values $p_e$ for surviving pairs are the same at both prefixes.
Since $c\ge2$, all $c-2p_e$ are nonnegative. Restriction cannot increase
$\operatorname{cost}_t(F)$, so the last sum and reference cost can both
be bounded by their values at $t$. Taking the maximum proves the result.
\end{proof}

\subsection{Exact costs and the original recourse guarantee}

\begin{theorem}\label{thm:improved-online-exact}
For every oblivious vertex-arrival instance with $m$ input partitions,
known horizon $n$, $\zeta\in(0,1)$, and
$0<\varepsilon<1/(\eta+1)$, there is a randomized online algorithm such
that, with probability at least $1-\zeta$,
\[
 \operatorname{cost}_t(C_t)
 \le \left(2-\frac{1}{\eta+1}+\varepsilon\right)
       \operatorname{OPT}_t
 \qquad\text{simultaneously for all }t\le n.
\]
On every execution, each vertex has recourse
\[
 O_\eta\!\left(
  \varepsilon^{-2}\log\frac{2n}{\zeta}\log\frac{2mn}{\zeta}
 \right).
\]
\end{theorem}

\paragraph{Algorithm.}
The cases $n=1$ and $m\le2$ use the optimal input-following rule in the
draft. In the remaining cases, use its candidate selection procedure with
\[
 w=\frac{\eta}{\eta+1},\qquad
 K=\left\lceil\frac4\varepsilon\log\frac{2n}{\zeta}\right\rceil,
 \qquad \theta=\frac\varepsilon4.
\]
Each candidate follows a uniform input with probability $w$, and otherwise
runs one persistent copy of Online Agreement on an independently sampled
mixed-majority graph. All candidate processes advance on every arrival.
For candidate $a$, compute its exact current cost $f_a(t)$ and history
$G_a(t)=\max_{u\le t}f_a(u)$, and put $M_t=\min_aG_a(t)$.
While $M_t=0$, follow a candidate with zero history, retaining the current
candidate whenever possible. A positive phase starts with $B=M_t$ and
threshold $(1+\theta)B$. Retain the selected candidate while its history
is below the threshold. If it becomes ineligible, select another eligible
candidate. If $M_t$ reaches the threshold, first begin a new phase with
baseline $M_t$ and select a candidate attaining this minimum. Break ties
by a fixed rule. Prefix the selected candidate's local identifiers by its
candidate index.

\begin{proof}
Fix a prefix $t$ and an optimal partition $F_t^*$. For every pair $e$,
let $p_e$ be the fraction of inputs disagreeing with $F_t^*$ on $e$, and
put $X_t=\sum_e p_e^2$. Then $\sum_e p_e=\operatorname{OPT}_t$.
A uniform input candidate has nondecreasing cost and satisfies
\[
 \mathbb E[G_a(t)\mid\text{input}]
 =2\operatorname{OPT}_t-2X_t.
\]
By Lemma~\ref{lem:improved-weighted-history} and then
Lemma~\ref{lem:improved-majority}, a mixed-majority candidate satisfies
\begin{align*}
 \mathbb E[G_a(t)\mid\text{mixed majority}]
 &\le\operatorname{OPT}_t+\sum_e(c-2p_e)h_{1/\eta}(p_e)\\
 &\le\operatorname{OPT}_t+2\eta X_t.
\end{align*}
The choice $w=\eta/(\eta+1)$ cancels $X_t$ exactly,
\[
 \mathbb E[G_a(t)]
 \le w(2\operatorname{OPT}_t-2X_t)
 +(1-w)(\operatorname{OPT}_t+2\eta X_t)
 =R_*\operatorname{OPT}_t.
\]
Let $R=R_*+\varepsilon/2<2$. If $\operatorname{OPT}_t>0$,
Markov's inequality gives
\[
 \Pr[G_a(t)\le R\operatorname{OPT}_t]
 \ge1-R_*/R=\varepsilon/(2R)\ge\varepsilon/4.
\]
If $\operatorname{OPT}_t=0$, all inputs coincide on the prefix and every
candidate has history zero. Independence of the candidates and the choice
of $K$ imply
\[
 \Pr[M_t>R\operatorname{OPT}_t]\le\zeta/(2n).
\]
A union bound gives this history bound simultaneously at every prefix.
The phase rule ensures
$\operatorname{cost}_t(C_t)\le(1+\theta)M_t$. On the event just obtained,
\[
 (1+\theta)R=R_*+\varepsilon/2+\varepsilon R/4
 <R_*+\varepsilon.
\]

For completeness, exact histories belong to
$m^{-1}\mathbb Z\cap[0,\binom n2]$. Positive phase baselines grow by at
least $1+\theta$, so there are $O(\varepsilon^{-1}\log(2mn))$
positive phases. Monotonicity of every history implies that, once a candidate becomes
ineligible, it remains ineligible until the next phase.
There are at most $K$ switches per positive phase and at most $K-1$ in
the zero stage. For any vertex, add these switches to the at most
$Kq_{\rm A}\lceil\log_2 n\rceil$ internal changes across all persistent
solver copies. This proves the asserted recourse for every execution.
\end{proof}

\paragraph{Optimality of this pairwise certificate.}
The constant $R_*$ is optimal within the following restricted analysis.
Mix a uniform input with probability $w$ and a sampled-majority solver
with probability $1-w$, use the weighted history comparison above, and
require a bound separately for every $p\in[0,1]$. For any fixed finite
mixture of odd majorities of size at least three, its error function
satisfies $h(p)/p\to0$ as $p\to0$ and $h(1/2)=1/2$.
The pairwise coefficient is
\[
 r_w(p)=w\,2(1-p)
       +(1-w)\left(1+(c-2p)\frac{h(p)}p\right).
\]
Its supremum is at least
\[
 \max\{\lim_{p\downarrow0}r_w(p),r_w(1/2)\}
 =\max\{1+w,\ 1+\eta(1-w)\}.
\]
Minimizing this maximum gives $w=\eta/(\eta+1)$ and $R_*$. This is a
limitation of the stated certificate. It gives no lower bound on the
online problem.

\subsection{Removing the input-count dependence from recourse}

The estimated-history construction in the explicit reduction of the draft
also applies to deterministic Online Agreement. It gives the following
stronger recourse bound while retaining the improved ratio.

\begin{theorem}\label{thm:improved-online-estimated}
Under the hypotheses of Theorem~\ref{thm:improved-online-exact}, one can
achieve the same simultaneous competitive ratio and failure probability
with per-vertex recourse
\[
 O_\eta\!\left(
   \varepsilon^{-2}\log\frac{2n}{\zeta}
                    \log\frac{2n}{\varepsilon\zeta}
 \right)
\]
on every execution. This bound is independent of $m$.
\end{theorem}

\paragraph{Algorithm.}
Use the same candidate distribution and set
\[
 \theta=\gamma=\varepsilon/40,\qquad
 K=\left\lceil\frac4\varepsilon\log\frac{4n}{\zeta}\right\rceil,
 \qquad L=\left\lceil\log_2\frac{4n}{\zeta}\right\rceil,
 \qquad g=\left\lceil\frac{2}{\gamma^2}
                   \log\frac{8Kn^2L}{\zeta}\right\rceil.
\]
Independently of the candidates, sample $L$ groups of $g$ independent
uniform input indices, with replacement. For a partition $C$ of $V_t$,
define
\[
 \widehat f_t^{(j)}(C)=\frac1g\sum_{b=1}^g d(C,P_{J_{j,b}}|_{V_t}),
 \qquad \widehat f_t(C)=\min_{j\in[L]}\widehat f_t^{(j)}(C).
\]
Use histories
$\widehat G_a(t)=\max_{u\le t}\widehat f_u(A_a(u))$ and
$\widehat M_t=\min_a\widehat G_a(t)$ in the phase rule.
Candidate processes themselves are independent of these evaluation samples.

\begin{proof}
Let $R=R_*+\varepsilon/2$. Define a proxy $Z_a(t)$ to be the true
cost for an input candidate and
$\operatorname{OPT}_t+\sum_{e\in B_{H_a}^t(F_t^*)}(c-2p_e)$ for a majority candidate.
The preceding expectation calculation gives
$\mathbb E Z_a(t)\le R_*\operatorname{OPT}_t$.
Thus, with probability at least $1-\zeta/4$, every prefix has a candidate
$b$ with $Z_b(t)\le R\operatorname{OPT}_t$.

With probability at least $1-\zeta/2$, the independent evaluation groups
simultaneously satisfy, for every $a$ and $u\le t$,
\begin{align}
 \widehat f_t(F_t^*)&\le2\operatorname{OPT}_t,
 \label{eq:improved-eval-ceiling}\\
 \left|\widehat f_u(A_a(u))-\widehat f_u(F_t^*|_{V_u})
       -\bigl(\operatorname{cost}_u(A_a(u))
              -\operatorname{cost}_u(F_t^*|_{V_u})\bigr)\right|
 &\le\gamma d(A_a(u),F_t^*|_{V_u}).
 \label{eq:improved-eval-comparison}
\end{align}
For the first inequality, each group mean exceeds
$2\operatorname{OPT}_t$ with probability at most $1/2$ by Markov's
inequality. Taking a minimum over the independent groups and then a union
bound over prefixes costs at most $n2^{-L}\le\zeta/4$.
For the second, condition on all candidate trajectories. Each variable
$d(A_a(u),P|_{V_u})-d(F_t^*|_{V_u},P|_{V_u})$ lies in
$[-d(A_a(u),F_t^*|_{V_u}),d(A_a(u),F_t^*|_{V_u})]$.
Hoeffding's inequality and a union bound over at most $Kn^2L$ comparisons
give failure at most
$2Kn^2L\exp(-g\gamma^2/2)\le\zeta/4$.
If two lists have all coordinatewise differences in a fixed interval,
the difference of their minima belongs to that interval. This proves
\eqref{eq:improved-eval-comparison}. The zero-distance cases are exact.

Fix a prefix where both events hold and a candidate $b$ with small proxy.
For an input candidate, every group cost is nondecreasing, so
$\widehat G_b(t)=\widehat f_t(A_b(t))$.
Its distance to $F_t^*$ is at most
$(R+1)\operatorname{OPT}_t<3\operatorname{OPT}_t$.
Equation~\eqref{eq:improved-eval-comparison} therefore implies
\[
 \widehat G_b(t)
 \le\widehat f_t(F_t^*)+(R-1+3\gamma)\operatorname{OPT}_t.
\]
For a majority candidate, let $F_u=F_t^*|_{V_u}$. The small proxy and
Lemma~\ref{lem:improved-weighted-history} give, for every $u\le t$,
\[
 \operatorname{cost}_u(A_b(u))-\operatorname{cost}_u(F_u)
 \le (R-1)\operatorname{OPT}_t.
\]
Also $\operatorname{cost}_u(F_u)\le\operatorname{OPT}_t$, so
\[
 d(A_b(u),F_u)
 \le\operatorname{cost}_u(A_b(u))+\operatorname{cost}_u(F_u)
 \le(R+1)\operatorname{OPT}_t<3\operatorname{OPT}_t.
\]
Equation~\eqref{eq:improved-eval-comparison} now implies
\[
 \widehat f_u(A_b(u))
 \le\widehat f_u(F_u)+(R-1+3\gamma)\operatorname{OPT}_t
 \le\widehat f_t(F_t^*)+(R-1+3\gamma)\operatorname{OPT}_t.
\]
The last inequality follows by restriction for every evaluation group
and then for their minimum. Take the maximum over $u$ to obtain the
same history bound as for the input candidate.
Thus this bound holds in both cases and also bounds
$\widehat M_t$.

The phase rule gives
$\widehat f_t(C_t)\le(1+\theta)\widehat M_t$.
Combining this with \eqref{eq:improved-eval-comparison}, the preceding
history bound, and
$d(C_t,F_t^*)\le\operatorname{cost}_t(C_t)+\operatorname{OPT}_t$
gives
\[
 (1-\gamma)\operatorname{cost}_t(C_t)
 \le\left(R+4\gamma+\theta(R+1+3\gamma)\right)
          \operatorname{OPT}_t.
\]
Here \eqref{eq:improved-eval-ceiling} bounds the term
$\theta\widehat f_t(F_t^*)$. Since $R<2$,
$\theta=\gamma=\varepsilon/40<1/80$, and
\[
 \frac{R+4\gamma+\gamma(R+1+3\gamma)}{1-\gamma}
 =R+\frac{\gamma(2R+5+3\gamma)}{1-\gamma}
 <R+10\gamma<R_*+\varepsilon,
\]
the approximation follows simultaneously at all prefixes, including
zero-optimum prefixes. The total failure probability is at most
$3\zeta/4$.

Every empirical score and its history belongs to
$g^{-1}\mathbb Z\cap[0,\binom n2]$. The exact-cost phase counting argument
now gives at most
\[
 K\left(2+\left\lceil\log_{1+\theta}(g n^2)\right\rceil\right)
 =O(K\varepsilon^{-1}\log(gn+2))
\]
outer switches on every execution. Add at most
$Kq_{\rm A}\lceil\log_2n\rceil$ internal changes for each vertex.
The parameter choices imply
$\log(gn+2)=O(\log(2n/(\varepsilon\zeta)))$, which gives the theorem.
\end{proof}

\subsection{A stronger recourse bound with high probability}

The factor $K$ in the switch bound can be replaced by $O(\log K)$
if the selector is randomized. This gives a guarantee in expectation and
with high probability, in addition to the preceding deterministic bound.

At the start of each positive phase, assign independent uniformly random
priorities to the candidates and select the lowest-priority eligible
candidate. Keep that candidate until it becomes ineligible, and then
choose the lowest-priority surviving eligible candidate. Use a fresh
independent ordering at the next phase. Use the same rule with one fresh
ordering in the initial zero-history stage. All these orderings are
independent of the candidates and evaluation samples. At the start of a
positive phase the chosen candidate can have history between $B$ and
$(1+\theta)B$. This still gives exactly the selection inequality used in
the approximation proof.

\begin{theorem}\label{thm:improved-random-recourse}
Use the estimated-history algorithm of
Theorem~\ref{thm:improved-online-estimated} with the randomized selector
above. Its approximation guarantee is unchanged. Write
\[
 J=2+\left\lceil\log_{1+\theta}(gn^2)\right\rceil,
 \qquad H_K=\sum_{j=1}^K\frac1j.
\]
For every $\delta\in(0,1)$, with probability at least $1-\delta$,
every vertex simultaneously has recourse at most
\[
 (e-1)JH_K+\log(1/\delta)
       +Kq_{\rm A}\lceil\log_2n\rceil.
\]
Its expected recourse, and the expectation of the maximum recourse over
vertices, are at most
$JH_K+Kq_{\rm A}\lceil\log_2n\rceil$.
In particular, the high-probability bound is
\[
 O_\eta\!\left(
  \varepsilon^{-1}\log\frac{2n}{\varepsilon\zeta}
       \log\left(2+\varepsilon^{-1}\log\frac{2n}{\zeta}\right)
 +\varepsilon^{-1}\log\frac{2n}{\zeta}\log(n+2)
 +\log\frac1\delta
 \right).
\]
The approximation and this recourse bound hold jointly with probability
at least $1-\zeta-\delta$. Applying the theorem with each failure
parameter equal to half of a desired total failure probability gives a
single simultaneous guarantee.
\end{theorem}
\begin{proof}
Condition on every candidate trajectory and evaluation sample. All
histories and phase boundaries are then fixed independently of the
selector's random choices. During a phase the eligible sets are nested.
Order the initially eligible candidates by the time at which they cease
to be eligible, breaking ties deterministically and placing candidates
that survive the phase at the end. Refine simultaneous removals into
individual removals. Each candidate selected during the actual phase
is also selected at some time during this refined removal process.
The number of selections in the refined process is the number of
right-to-left record minima in a uniform random permutation of the
$r\le K$ candidate priorities.

Let $N_r$ denote this number. Its distribution is a sum of independent
Bernoulli variables with success probabilities $1,1/2,\ldots,1/r$.
For completeness, the relative rank of the $j$th element of a uniform
random permutation among its first $j$ elements is independent and
uniform on $[j]$. The new element is a record precisely when this rank
is one. Reading the permutation in reverse gives the stated distribution.
Consequently,
\[
 \mathbb E N_r=H_r\le H_K,
 \qquad
 \mathbb E e^{N_r}
 =\prod_{j=1}^r\left(1+\frac{e-1}{j}\right)
 \le\exp((e-1)H_K).
\]
There are at most $J$ stages including the zero-history stage, and their
random orderings are independent conditional on the trajectories.
Writing $S$ for the number of output switches, the total number of
selections bounds $S$, so
\[
 \mathbb E[S\mid\text{trajectories}]\le JH_K,
 \qquad
 \mathbb E[e^S\mid\text{trajectories}]
 \le\exp((e-1)JH_K).
\]
Markov's inequality gives
\[
 \Pr[S>(e-1)JH_K+\log(1/\delta)
        \mid\text{trajectories}]\le\delta.
\]
These inequalities hold for every conditioning and therefore hold
unconditionally. The same switch count applies to every vertex. Adding
the deterministic bound on internal identifier changes proves the
simultaneous statement, the bound on the expected maximum, and the
asymptotic expression.
\end{proof}

\subsection{Conditional transfer to an explicit dynamic solver}
\label{sec:conditional}

Suppose an explicit correlation-clustering solver has the hypothesis
required by the draft's consensus reduction, namely simultaneous
$\eta$-approximation with probability at least $1/2$ and expected prefix
work at most $tT_{\mathcal D}(n)$ with $T_{\mathcal D}(n)=\operatorname{polylog}(n+2)$.
The improved mixed-majority construction transfers to this setting.
Use the preceding evaluation groups and double $K$ to
$\lceil8\varepsilon^{-1}\log(4n/\zeta)\rceil$.
Conditional on any sampled majority graph, its solver succeeds throughout
the stream with probability at least $1/2$. Thus a candidate has a small
proxy and successful solver with probability at least $\varepsilon/8$,
and the same existence and estimated-history proof applies. The resulting
competitive ratio is $R_*+\varepsilon$, with the update-time bound
\[
 O\!\left(
  \varepsilon^{-3}\log^2\frac{2n}{\zeta}
  \log\frac{2n}{\varepsilon\zeta}\cdot\operatorname{polylog}(n+2)
 \right).
\]
This statement is conditional on the solver hypothesis. It does not
supply the missing solver guarantees identified in the review.

The graph queries needed by the reduction can be implemented explicitly.
For $k\in\{5,15\}$ and $q=(k+1)/2$, the positive closed neighborhood is
the union of the $\binom{k}{q}$ buckets formed by agreeing with the query
vertex on a specified $q$-subset of sampled coordinates. To support exact
uniform sampling with a worst-case bound, maintain buckets for all
subsets $S\subseteq[k]$ of size at least $q$, ordered by vertex arrival
index. There are 16 such buckets per vertex for $k=5$ and 16384 for $k=15$.
These counts are absolute constants, although the latter substantially
increases the constants in the update time and space.
Let $B_S(v)$ denote the bucket containing $v$ and let $F_v(x)$ count its
closed positive neighbors of arrival index at most $x$. Then
\[
 F_v(x)=\sum_{r=q}^k(-1)^{r-q}\binom{r-1}{q-1}
              \sum_{\substack{S\subseteq[k]\\|S|=r}}
                    |B_S(v)\cap[ x ]|.
\]
Indeed, a vertex agreeing in $d$ coordinates receives coefficient
\[
 \sum_{r=q}^d(-1)^{r-q}\binom{r-1}{q-1}\binom dr
 =\mathbf1[d\ge q].
\]
For $d=q$ the sum is one, and its increment from $d$ to $d+1$ is
$\binom d{q-1}(1-1)^{d-q+1}=0$ for $d\ge q$.
Order-statistic trees evaluate $F_v(x)$ in $O(\log(n+2))$ time.
The open-neighborhood prefix count is
$F_v(x)-\mathbf1[\operatorname{index}(v)\le x]$.
Choose a uniform integer rank in the open neighborhood and binary search
$[n]$ for its vertex using these prefix counts. This takes
$O(\log^2(n+2))$ worst-case time and samples exactly uniformly.
Insertions and degree queries take $O(\log(n+2))$, and edge queries use
$O(1)$ label comparisons. The extra logarithm in sampling is absorbed in
the displayed polylogarithmic update-time factor.

\section{A \texorpdfstring{$1.64+\varepsilon$}{1.64 + epsilon} approximation with the original seven samples}
\label{sec:centered-finite}

The approximation ratio can be improved without enlarging the sample.
The seven-input majority used in the draft already gives
$41/25+\varepsilon=1.64+\varepsilon$, with its original $35$ buckets per
vertex. There is also a $5/3+\varepsilon$ implementation using only five
sampled inputs and $\binom53=10$ buckets per vertex. Both conclusions
hold for the offline and implicit fully dynamic algorithms, with the
same candidate generation, repetition, evaluation, and maintenance
procedures. The improvement comes from a sampling correction whose
expectation is exactly zero, rather than an upper bound on the absolute
cost of a sampled sign error.

For odd $k$, define
\[
 \phi_k(s)=\sum_{j=(k+1)/2}^k\binom kj s^j(1-s)^{k-j},
 \qquad t_k(s)=(1-2s)\phi_k(s),\quad 0\le s\le1.
\]
For a signed graph $H$, let $s_e$ be the consensus cost of following its
sign on pair $e$. Thus the opposite sign has cost $1-s_e$, and the
expected contribution of a uniformly sampled input partition is
$z_e=2s_e(1-s_e)$.

\begin{lemma}[A correction with zero expectation]
\label{lem:centered-zero}
If $H$ is the majority of $k$ independently sampled uniform input
partitions, then $\mathbb E_H[t_k(s_e)]=0$ for every pair $e$.
\end{lemma}
\begin{proof}
Write $m_e$ for the smaller of the two possible consensus costs on
$e$, and $p=\phi_k(m_e)$. Except at a tie, where $t_k(s_e)=0$,
$s_e=m_e$ with probability $1-p$ and $s_e=1-m_e$ with probability $p$.
Oddness of $k$ gives $\phi_k(1-m_e)=1-p$, so
\[
 \mathbb E_H[t_k(s_e)]
 =(1-p)(1-2m_e)p+p(2m_e-1)(1-p)=0.
\]
No independence between different edges is required.
\end{proof}

\begin{theorem}\label{thm:centered-finite}
Let $C_H$ be a uniformly randomized Pivot candidate on the majority
of $k$ independent uniform input samples. Let $P_J$ be an independently
sampled uniform input partition. Either row of the following table
satisfies
\[
 \beta\,\mathbb E[\operatorname{cost}(C_H)]
 +(1-\beta)\,\mathbb E[\operatorname{cost}(P_J)]
 \le\gamma\operatorname{OPT}.
\]
\[
\begin{array}{c|c|c|c|c}
 k&\beta&\lambda&\gamma&K_{\rm buck}\\\hline
 5&1/3&4/3&5/3&10\\
 7&9/25&29/20&41/25&35
\end{array}
\]
Consequently the offline and implicit fully dynamic frameworks attain
$\gamma+\varepsilon$ with probability at least $1-\zeta$, using the
draft's asymptotic time and space bounds. In the dynamic setting this
holds simultaneously at every update of the prescribed oblivious
sequence.
\end{theorem}
\begin{proof}
Fix one row of the table and put
\begin{align*}
 f(s)&=\beta(1-s)+(1-\beta)2s(1-s)+\lambda t_k(s),\\
 g(s)&=\beta s+(1-\beta)2s(1-s)+\lambda t_k(s)
      =f(s)-\beta(1-2s).
\end{align*}
We will establish the local inequalities
\begin{align}
 g(s)&\le\gamma\min(s,1-s),\quad 0\le s\le1,
   \label{eq:finite-centered-pair}\\
 f(s_1)+f(s_2)+f(s_3)&\le\gamma C_F
   \label{eq:finite-centered-triplet}
\end{align}
for every bad triangle of $H$ and every reference partition $F$.
Number its positive edges $1,2$ and its negative edge $3$. The triangle
inequality for consensus distances gives $s_1+s_2+s_3\ge1$.
The five possible local reference costs are
\[
\begin{split}
 C_1&=s_1+s_2+1-s_3,\qquad
 C_2=s_1+1-s_2+s_3,\qquad
 C_3=1-s_1+s_2+s_3,\\
 C_4&=3-s_1-s_2-s_3,\qquad
 C_5=2-s_1-s_2+s_3.
\end{split}
\]

Here are exact polynomial certificates for all the needed inequalities.
The two polynomials $f$ are
\begin{align*}
 f_5(s)&=\frac13+s-\frac43s^2+\frac{40}{3}s^3
             -\frac{140}{3}s^4+48s^5-16s^6,\\
 f_7(s)&=\frac9{25}+\frac{23}{25}s-\frac{32}{25}s^2
             +\frac{203}{4}s^4-\frac{2233}{10}s^5
             +\frac{3451}{10}s^6-232s^7+58s^8.
\end{align*}
For a polynomial $p(s)$ on $[a,b]$, substitute $s=a+(b-a)x$.
If the resulting power coefficients are $A_i$, its degree-$d$
Bernstein coefficients are
\[
 B_j=\sum_{i=0}^j A_i\frac{\binom ji}{\binom di},
 \qquad j=0,\ldots,d.
\]
The polynomial is at least $\min_j B_j$ throughout the interval.
Define $q(s)=\gamma(1-s)-f(s)$. Direct rational arithmetic gives
\[
\begin{array}{c|c|c|c}
 p(s)&[a,b]&\min B_j\ (k=5)&\min B_j\ (k=7)\\\hline
 {}[\gamma s-g(s)]/s^2 &[0,1/2]&\text{see below}&221/1600\\
 \gamma(1-s)-g(s)&[1/2,1]&5/18&457/1600\\
 \gamma-f'(s)&[0,1/2]&1/5&761/1400\\
 -\gamma-f'(s)&[1/2,1]&0&17/100\\
 q(s)&[0,1/2]&1/3&99/320
\end{array}
\]
Use the actual degree of each displayed polynomial: the degrees in
these five rows are $(k-1,k+1,k,k,k+1)$. The first quotient extends
polynomially at zero. For $k=5$, its Bernstein minima on $[0,1/4]$
and $[1/4,1/2]$ are respectively $5/48$ and $11/48$.
Thus \eqref{eq:finite-centered-pair} follows. These certificates also prove
\[
 f'(s)\le\gamma\quad(0\le s\le1),\qquad
 f'(s)\le-\gamma\quad(1/2\le s\le1),\qquad
 q(s)\ge0\quad(0\le s\le1/2).
\]
Since $q'=-\gamma-f'\ge0$ on $[1/2,1]$ and
$f(1/2)=1/2$, they further give
\begin{equation}
 q(s)\ge(\gamma-1)/2>0\qquad(1/2\le s\le1).
 \label{eq:centered-q}
\end{equation}

Consider first $C_1$. The gap $\gamma C_1-\sum_i f(s_i)$ is
nondecreasing in $s_1,s_2$. Under $s_1+s_2+s_3\ge1$, it is enough
to check the boundary $s_1+s_2=1-s_3$: decrease $s_1,s_2$ to that
boundary while keeping them nonnegative. The certificates below prove
the inequality on this enlarged domain; no other feasibility
constraints are needed. Put $t=s_1+s_2$ and, when $t>0$,
$w=(s_1-s_2)^2/t^2$. Then $0\le t,w\le1$ and the required gap is
\[
 P_k(t,w)=2\gamma t-f(1-t)
     -f\left(\frac{t(1+\sqrt w)}2\right)
     -f\left(\frac{t(1-\sqrt w)}2\right).
\]
Odd powers of $\sqrt w$ cancel, so this is a polynomial of bidegree
$(k+1,(k+1)/2)$. If $f(s)=\sum_{i=0}^{k+1}a_i s^i$, an explicit
formula without radicals is
\begin{equation}
 P_k(t,w)=2\gamma t-\sum_{i=0}^{k+1}a_i(1-t)^i
 -\sum_{i=0}^{k+1}\sum_{j=0}^{\lfloor i/2\rfloor}
      2^{1-i}a_i\binom i{2j}t^i w^j.
 \label{eq:centered-P}
\end{equation}
The formula also handles $t=0$.

For completeness, the following is an exact tensor Bernstein
certificate for $P_k\ge0$ on the entire unit square. On a rectangle
$[a,b]\times[c,d]$, substitute $t=a+(b-a)x$ and $w=c+(d-c)y$.
If the resulting power coefficients are $A_{ij}$, the tensor
Bernstein coefficients at bidegree $(D,E)=(k+1,(k+1)/2)$ are
\[
 B_{rs}=\sum_{i=0}^{r}\sum_{j=0}^{s}
 A_{ij}\frac{\binom ri}{\binom Di}\frac{\binom sj}{\binom Ej}.
\]
Each value of the polynomial is a convex combination of these
coefficients. The following nine rectangles cover $[0,1]^2$ and have
strictly positive coefficient minima:
\[
\begin{array}{c|c|c|c}
 {}[a,b]&[c,d]&\min B_{rs}\ (k=5)&\min B_{rs}\ (k=7)\\\hline
 {}[0,1/4]&[0,1/2]&6213/32768&200354911/838860800\\
 {}[0,1/4]&[1/2,1]&37433/196608&80256847/335544320\\
 {}[1/4,3/8]&[0,1/4]&39163/1572864&4924796569/107374182400\\
 {}[1/4,1/2]&[1/4,1/2]&1473/163840&14705001/2936012800\\
 {}[1/4,1/2]&[1/2,1]&4169/122880&248019/10485760\\
 {}[3/8,7/16]&[0,1/8]&6253/16777216&5077941423/1374389534720\\
 {}[3/8,1/2]&[1/8,1/4]&1872017/150994944&19395179149/3006477107200\\
 {}[7/16,1/2]&[0,1/8]&15889/33554432&260187801/375809638400\\
 {}[1/2,1]&[0,1]&11/384&7247/409600
\end{array}
\]
All entries are obtained directly from \eqref{eq:centered-P} and
the displayed coefficient conversion. This proves the $C_1$ case;
$C_2,C_3$ follow by permuting coordinates.

For $C_4$, the gap is $q(s_1)+q(s_2)+q(s_3)\ge0$.
For $C_5$, if $s_1\le1/2$, then $C_5\ge C_2$; if $s_2\le1/2$,
then $C_5\ge C_3$. Otherwise $s_1,s_2\ge1/2$. Its gap is
nondecreasing in $s_3$, since its derivative is $\gamma-f'(s_3)$.
At $s_3=0$, \eqref{eq:centered-q} bounds the gap below by
\[
 q(s_1)+q(s_2)-f(0)\ge\gamma-1-\beta=1-2\beta>0,
\]
where both rows satisfy $\gamma=2-\beta$. This finishes the proof
of \eqref{eq:finite-centered-triplet}.

To sum these local inequalities, fix $H$ and a reference partition $F$.
Let $D_e$ be Pivot's probability of opposing the sign of $e$. The
exact Pivot identity from the draft is
$D_e=\sum_{T\ni e}p_T/3$, where the sum is over bad triangles of $H$.
Multiply \eqref{eq:finite-centered-pair} for $e$ by $1-D_e$ and
\eqref{eq:finite-centered-triplet} for $T$ by $p_T/3$, and sum.
Each reference edge cost, uniform-input contribution, and correction
$\lambda t_k(s_e)$ receives total weight one. Consequently
\[
 \beta\,\mathbb E_{\rm piv}[\operatorname{cost}(C_H)]
 +(1-\beta)\,\mathbb E_J[\operatorname{cost}(P_J)]
 +\lambda\sum_e t_k(s_e)
 \le\gamma\operatorname{cost}(F).
\]
Taking expectation over $H$ and applying Lemma~\ref{lem:centered-zero}
cancels the last sum exactly. Choose $F$ optimal to obtain the theorem's
expected-cost bound. The expected minimum of the two candidates is
no larger than this convex combination. The draft's independent
repetition and evaluation argument gives $\gamma+\varepsilon$ with
failure probability $\zeta$; in the dynamic setting use the same
fixed-update distribution and union bound as in the draft. All sample
sizes here are fixed constants, giving the asserted time bounds.
\end{proof}

\subsection{Further finite-sample tradeoffs}

The same proof gives three additional approximation factors, including
$8/5+\varepsilon=1.6+\varepsilon$ with fifteen sampled inputs. The
following rows can be added to Theorem~\ref{thm:centered-finite}:
\[
\begin{array}{c|c|c|c|c}
 k&\beta&\lambda&\gamma&K_{\rm buck}\\\hline
 9&3/8&8/5&13/8&126\\
 13&39/100&9/5&161/100&1716\\
 15&2/5&9/5&8/5&6435
\end{array}
\]
In every row $\gamma=2-\beta$, and $f,g,q,P_k$ retain exactly their
preceding definitions. Here are all the additional univariate
certificates needed in that proof. The degree conventions are unchanged.
\[
\begin{array}{c|c|c|c|c}
 p(s)&[a,b]&\min B_j\ (k=9)&\min B_j\ (k=13)&\min B_j\ (k=15)\\\hline
 {}[\gamma s-g(s)]/s^2 &[0,1/2]&89/280&101/200&15927/28160\\
 \gamma(1-s)-g(s)&[1/2,1]&117/400&417/1400&47/160\\
 \gamma-f'(s)&[0,1/2]&49/90&5599/8320&95/128\\
 -\gamma-f'(s)&[1/2,1]&7/20&29/50&3/5\\
 q(s)&[0,1/2]&59/200&199/700&9/32
\end{array}
\]
For each $k$, the following rectangles cover the unit square and give
positive tensor Bernstein coefficients for $P_k$, at the same
bidegree $(k+1,(k+1)/2)$. They complete the certificate for the
triplet inequality, so no new argument is required.
\[
\begin{array}{c|c|c|c}
 k&[a,b]&[c,d]&\min B_{rs}\\\hline
 9 &[0,1/4]&[0,1/2]&25970139/83886080\\
 9 &[0,1/2]&[1/2,1]&13267/14745600\\
 9 &[1/4,3/8]&[0,1/4]&8242846289791/109951162777600\\
 9 &[1/4,1/2]&[1/4,1/2]&76076627/15099494400\\
 9 &[3/8,1/2]&[0,1/4]&1249243/314572800\\
 9 &[1/2,1]&[0,1]&1267/81920\\\hline
 13 &[0,1/4]&[0,1/2]&2792808933987/6871947673600\\
 13 &[0,1/2]&[1/2,1]&8763372747/610690662400\\
 13 &[1/4,1/2]&[0,1/2]&57195679509/8549669273600\\
 13 &[1/2,1]&[0,1]&3526147/293601280\\\hline
 15 &[0,1/4]&[0,1/2]&17437117345851/43980465111040\\
 15 &[0,1/2]&[1/2,1]&5669651969/858993459200\\
 15 &[1/4,1/2]&[0,1/4]&4082432264899/3078632557772800\\
 15 &[1/4,1/2]&[1/4,1/2]&1832498068713/879609302220800\\
 15 &[1/2,1]&[0,1]&8259643/1342177280
\end{array}
\]


\subsection{A fixed factor without an accuracy parameter}
\label{sec:fixed-ratio-improvement}

The nine-sample bound $13/8$ gives a fixed approximation factor $41/25=1.64$
with no accuracy parameter in the running time. This uses 126 buckets
per vertex. The seven-sample result above retains 35 buckets and allows
an arbitrary additive $\varepsilon$ instead.

Here is an explicit amplification argument, also recording the
evaluation guarantee used by the finite and asymptotic results.
Set $\gamma_0=13/8$ and let $H_0=1$ for one offline instance, or
$H_0=\max\{1,U\}$ for a dynamic stream of at most $U$ updates. Choose
\[
 \delta=\frac1{200},\qquad
 T=\left\lceil400\log\frac{2H_0}{\zeta}\right\rceil,
 \qquad
 s=\left\lceil720000\log\frac{4TH_0}{\zeta}\right\rceil.
\]
Generate both candidates in each of $T$ independent trials, use a common
independent sample of $s$ input partitions to estimate all their costs,
and return a minimum empirical-cost candidate.

For an instance with positive optimum, Markov's inequality gives a
trial whose better candidate costs at most
$(\gamma_0+\delta)\operatorname{OPT}$ with probability at least
$\delta/(\gamma_0+\delta)\ge\delta/2=1/400$.
The probability that any of the $H_0$ instances lacks such a trial is
at most $\zeta/2$. For zero optimum, all inputs coincide and both types
of candidates are exact.

Condition on all candidate trajectories and choose a true minimum-cost
candidate $C_0$ at each instance. Consider any candidate $C$ with
$\operatorname{cost}(C)>(1+\delta)\operatorname{cost}(C_0)$ and put
$R=d(C,C_0)$. An evaluation difference
$d(C,P)-d(C_0,P)$ lies in $[-R,R]$. Its mean is greater than
$\delta R/(2+\delta)$, since
$R\le\operatorname{cost}(C)+\operatorname{cost}(C_0)$.
Hoeffding's inequality bounds the probability that $C$ has empirical
cost no greater than $C_0$ by
$\exp(-s\delta^2/18)$, using $\delta\le1$.
A union bound over at most $2TH_0$ comparisons costs at most $\zeta/2$.
Thus, with probability at least $1-\zeta$, the chosen output satisfies
\[
 \operatorname{cost}(C)
 \le(1+\delta)(\gamma_0+\delta)\operatorname{OPT}
 =\frac{32763}{20000}\operatorname{OPT}
 <\frac{41}{25}\operatorname{OPT}.
\]
The statement holds simultaneously at every update in the dynamic
setting. The displayed sample constants are large; they are explicit
sufficient bounds, not optimized implementation parameters.

Since $\delta$ and the bucket count are fixed, the offline running time is
\[
 O\!\left(n\log^2\frac{n+2}{\zeta}\right).
\]
With the hashing implementation of Section~\ref{sec:hashing}, the
expected amortized dynamic update time is
\[
 O\!\left(\log^2\frac{U+2}{\zeta}\,\log(N+2)\right),
\]
with worst-case $O(1)$ membership and same-cluster queries. More
generally, the same fixed-slack argument applies to any fixed target
strictly larger than a certified expected factor.

\section{Polynomial dependence on the accuracy near \texorpdfstring{$11/7$}{11/7}}
\label{sec:centered-poly}

\subsection{The \texorpdfstring{$11/7+\varepsilon$}{11/7 + epsilon} guarantees}

\begin{theorem}[Offline and implicit fully dynamic consensus clustering]
\label{thm:eleven-sevenths}
Assume the draft's constant-time access to input cluster labels,
uniform sampling of input indices, and the word-RAM conventions of
Section~\ref{sec:hashing}. For every
$0<\varepsilon\le1$ and $0<\zeta<1$, there are randomized algorithms
with the following guarantees.
\begin{enumerate}
\item Given $m$ input partitions of $n$ vertices, the offline algorithm
returns a partition $C$ such that, with probability at least $1-\zeta$,
\[
 \operatorname{cost}(C)
 \le\left(\frac{11}{7}+\varepsilon\right)\operatorname{OPT}.
\]
Its running time is
\[
 n\varepsilon^{-O(1)}
       \log^2\frac{n+2}{\varepsilon\zeta}.
\]
\item Starting from the empty vertex set, the implicit fully dynamic
algorithm supports every oblivious sequence of at most a known number
$U\ge1$ of vertex insertions and deletions. With probability at least
$1-\zeta$, its maintained partitions satisfy
\[
 \operatorname{cost}_t(C_t)
 \le\left(\frac{11}{7}+\varepsilon\right)\operatorname{OPT}_t
 \qquad\text{simultaneously after every update}.
\]
Writing $N$ for the maximum active size, its expected amortized update
time is
\[
 \varepsilon^{-O(1)}
       \log^2\frac{U+2}{\varepsilon\zeta}\,\log(N+2).
\]
Membership and same-cluster queries take $O(1)$ worst-case time.
With $n_t$ active vertices, listing all vertex-label pairs takes
$O(n_t+1)$ time, and the maintained data use
\[
 (n_t+1)\varepsilon^{-O(1)}
       \log^2\frac{U+2}{\varepsilon\zeta}
\]
words of space, in addition to access to the input labels.
\end{enumerate}
The exponents hidden in $O(1)$ are absolute constants. The sample and
bucket counts depend on $\varepsilon$, with a minimum sample size of
87 in the certificate below.
\end{theorem}

\begin{proof}[Derivation from the single-trial bound]
The charging argument below proves \eqref{eq:centered-main}.
Set $\varepsilon_0=\min\{\varepsilon,1/2\}$ and choose the smallest
odd integer $k$ satisfying
\[
 k\ge87,\qquad k\ge10\log\frac{144}{7\varepsilon_0}.
\]
Then one trial's better candidate has expected cost at most
$\gamma_0\operatorname{OPT}$, where
$\gamma_0=11/7+\varepsilon_0/2$.
Let $H_0=1$ for the offline case and $H_0=U$ for the dynamic case.
Choose independent trials and an independent common evaluation sample
with parameters
\[
 \delta=\frac{\varepsilon_0}{8},\qquad
 T=\left\lceil\frac{2}{\delta}\log\frac{2H_0}{\zeta}\right\rceil,
 \qquad
 s=\left\lceil\frac{18}{\delta^2}
                  \log\frac{4TH_0}{\zeta}\right\rceil.
\]
At each instance or update, generate or maintain both candidates of
every trial and select a candidate of minimum empirical cost.
In the dynamic case, retain the samples throughout the stream and
maintain the trial processes with pivot priorities independent of the
evaluation sample, as in the draft.
The amplification and evaluation argument in
Section~\ref{sec:fixed-ratio-improvement} applies with these parameters.
In particular, $\gamma_0+\delta<2$, so a trial has cost at most
$(\gamma_0+\delta)\operatorname{OPT}$ with probability at least
$\delta/2$. Independent repetition and a union bound over the $H_0$
instances bound the probability of lacking such a candidate by
$\zeta/2$. Conditional on all candidate trajectories, the evaluation
sample selects a candidate within a factor $1+\delta$ of the best
true candidate cost at all $H_0$ instances, except with probability
$\zeta/2$. Consequently the output factor is at most
\[
 (1+\delta)(\gamma_0+\delta)
 =\frac{11}{7}+\frac{23}{28}\varepsilon_0
                 +\frac{5}{64}\varepsilon_0^2
 <\frac{11}{7}+\varepsilon_0
 \le\frac{11}{7}+\varepsilon.
\]
Zero-optimum instances are exact, as all input partitions coincide.
The dynamic conclusion uses the exact maintenance of each trial and
the fixed-update distribution on an oblivious stream.

Here $k=O(\log(2/\varepsilon))$ and
$K_{\rm buck}=\binom{k}{(k+1)/2}\le2^k=\varepsilon^{-O(1)}$.
The draft's bucket implementation has polynomial dependence on $k$
and $K_{\rm buck}$. Combining this dependence with the values of
$T,s$ and the hashing implementation in Section~\ref{sec:hashing}
gives the stated running-time, space, and query bounds. The dynamic
time bound is an unconditional expectation over the algorithm's
randomness.
\end{proof}

\subsection{The single-trial estimate}

The earlier exponential dependence on accuracy can be replaced by a polynomial. The new ingredient
is a signed correction whose expectation is exactly zero. This makes it
possible to retain the cancellation between the correct and incorrect
sampled signs, rather than paying for the incorrect signs alone.

Use the scaled correction $T_k=18t_k$ from Section~\ref{sec:centered-finite}.
For odd $k$, define
\[
 \phi_k(s)=\sum_{j=(k+1)/2}^k\binom{k}{j}s^j(1-s)^{k-j},
 \qquad T_k(s)=18(1-2s)\phi_k(s).
\]
For every odd $k\ge87$, we prove the single-trial bound
\begin{equation}
 \frac37\mathbb E[\operatorname{cost}(C_H)]
 +\frac47\mathbb E[\operatorname{cost}(P_J)]
 \le \left(\frac{11}{7}+\frac{72}{7}e^{-k/10}\right)
              \operatorname{OPT},
 \label{eq:centered-main}
\end{equation}
where $H$ is the majority graph of $k$ independent uniformly sampled
input partitions, $C_H$ is Pivot on that graph, and $P_J$ is an independent
uniform input. In particular, $k=O(\log(1/\tau))$ suffices for expected
factor $11/7+\tau$. The clique-bucket count is then polynomial in
$1/\tau$, improving the earlier exponential dependence.

\subsection{Univariate inequalities and exact certificates}

The following inequalities hold for every odd $k\ge87$:
\begin{align}
 T_k(s)&\le8s^2 &&(0\le s\le1/2),\label{eq:centered-pair-bound}\\
 T_k(s)&\le8(1-s)^2 &&(0\le s\le1),\label{eq:centered-allflip-bound}\\
 16s-8s^2+T_k(s)&\le13/2 &&(0\le s\le1),
       \label{eq:centered-singleton-bound}\\
 T_k(s)&\le8(s-1/4)^2+18e^{-k/10} &&(0\le s\le1).
       \label{eq:centered-critical-bound}
\end{align}
Here are finite exact certificates and their extension to all such $k$.
If $k=2h+1$, a binomial calculation gives
\[
 \phi_{k+2}(s)-\phi_k(s)
 =\binom{k}{h}s^{h+1}(1-s)^{h+1}(2s-1).
\]
Consequently $T_{k+2}(s)\le T_k(s)$ on the entire interval $[0,1]$.
Thus the first three inequalities follow from the corresponding
certificates for $k=13$ below. For a polynomial
$p(s)=\sum_{i=0}^d a_i s^i$ on an interval $[l,h]$, put
\[
 c_i=(h-l)^i\sum_{j=i}^d a_j\binom ji l^{j-i},\qquad
 b_j=\sum_{i=0}^j c_i\frac{\binom ji}{\binom di}.
\]
Then $b_j$ are its Bernstein coefficients after mapping $[l,h]$ to
$[0,1]$. The coefficients of $\phi_k$ used in this calculation are
\[
 {}[s^j]\phi_k(s)=
 \begin{cases}
 0,&j<(k+1)/2,\\
 (-1)^{j-(k+1)/2}\binom kj\binom{j-1}{(k-1)/2},&j\ge(k+1)/2.
 \end{cases}
\]
All computations in the following table are rational. In the rows
listing two intervals, the bound holds for every coefficient on both
intervals. The first quotient extends polynomially at zero.
\[
\begin{array}{c|c|c|c}
 p(s)&d&[l,h]&\hbox{lower bound on every }b_j\\\hline
 8-T_{13}(s)/s^2&12&[0,1/2]&17/20\\
 8(1-s)^2-T_{13}(s)&14&[0,1/2],\ [1/2,1]&21269/13312\\
 13/2-16s+8s^2-T_{13}(s)&14&[0,1/2],\ [1/2,1]&1301/13312\\
 8(s-1/4)^2-T_{87}(s)&88&[17/64,49/128],\ [49/128,1/2]&1/1000
\end{array}
\]
The first three rows prove
\eqref{eq:centered-pair-bound}--\eqref{eq:centered-singleton-bound}.
The last row and monotonicity in $k$ prove
$T_k(s)\le8(s-1/4)^2$ for $17/64\le s\le1/2$.
For $s\ge1/2$, $T_k(s)\le0$. For $0\le s\le17/64$, Hoeffding's
inequality gives
\[
 T_k(s)\le18\phi_k(s)
 \le18\exp\bigl(-2k(1/2-s)^2\bigr)
 \le18e^{-225k/2048}\le18e^{-k/10}.
\]
These three cases prove \eqref{eq:centered-critical-bound}.

\subsection{Local charging with the centered correction}

Fix an arbitrary complete signed graph $H$. For each pair $e$, let
$s_e$ be the consensus cost of following its sign in $H$,
$z_e=2s_e(1-s_e)$, and $c_e$ its cost in a reference partition $F$.
The corrected pair inequality is
\begin{equation}
 3s_e+4z_e+T_k(s_e)\le11c_e.
 \label{eq:centered-pair}
\end{equation}
It suffices to consider $c_e=\min(s_e,1-s_e)$. If $s_e\le1/2$,
the left side minus $11c_e$ is
$-8s_e^2+T_k(s_e)\le0$. If $s_e\ge1/2$, use
$\phi_k(s_e)\ge1/2$ and hence $T_k(s_e)\le9-18s_e$.
The difference is then at most $4s_e-8s_e^2-2\le0$.

For a bad triplet of $H$, number its positive edges 1 and 2 and its
negative edge 3, and put $S=s_1+s_2+s_3$. Consensus triangle
inequalities give $S\ge1$. We will show, for every reference
partition's local cost $C_F$, that
\begin{equation}
 \sum_{i=1}^3\{3(1-s_i)+4z_i+T_k(s_i)\}
 \le(11+72e^{-k/10})C_F.
 \label{eq:centered-triplet}
\end{equation}
Subtract $11C_F$ from the left side and call the result $G$.
There are five possible reference partitions.

First consider $C_F=a+b+1-r$, where $(a,b,r)$ is a permutation of
$(s_1,s_2,s_3)$. These are the three partitions that reverse exactly
one of the bad triplet's signs. Write $d=a+b+r-1\ge0$.
For every $r\in[0,1]$ we have
$T_k(r)\le9-18r$: below $1/2$ use $\phi_k(r)\le1/2$, and above
$1/2$ reverse the inequality when multiplying by $1-2r$.
Expanding around $(a,b,r)=(1/4,1/4,1/2)$ gives
\begin{align*}
 G
 &=-2-6a-6b+16r-8(a^2+b^2+r^2)
                  +T_k(a)+T_k(b)+T_k(r)\\
 &\le -10d-8\bigl((a-1/4)^2+(b-1/4)^2+(r-1/2)^2\bigr)
                  +T_k(a)+T_k(b)\\
 &\le36e^{-k/10}-8(r-1/2)^2-10d.
\end{align*}
If $r\le3/4$, then $C_F\ge2(1-r)\ge1/2$, and therefore
$G\le72e^{-k/10}C_F$. If $r>3/4$, the last upper bound is
strictly less than $36e^{-k/10}-1/2<0$ for $k\ge87$.

For the reference partition reversing all three signs,
$C_F=3-S$. Its gap is
\[
 G=\sum_{i=1}^3\{T_k(s_i)-8(1-s_i)^2\}\le0
\]
by \eqref{eq:centered-allflip-bound}.
For the all-singletons partition, $C_F=2-s_1-s_2+s_3$, giving
\[
 G=-13+\sum_{i=1}^2\{16s_i-8s_i^2+T_k(s_i)\}
                  -6s_3-8s_3^2+T_k(s_3).
\]
Each braced term is at most $13/2$ by
\eqref{eq:centered-singleton-bound}. The remaining term is nonpositive:
for $s_3\le1/2$ use \eqref{eq:centered-pair-bound}, and for
$s_3\ge1/2$ use $T_k(s_3)\le0$. This proves
\eqref{eq:centered-triplet} in all cases.

\subsection{Cancellation in expectation and consequences}

Multiply \eqref{eq:centered-pair} by $1-D_e$, where $D_e$ is Pivot's
probability of opposing pair $e$, and multiply each corrected triplet
inequality by $p_t/3$ as in the exact pair identity. Summing gives
\[
 3\mathbb E_{\rm piv}[\operatorname{cost}(C_H)]
 +4\mathbb E_J[\operatorname{cost}(P_J)]
 +\sum_e T_k(s_e)
 \le(11+72e^{-k/10})\operatorname{cost}(F).
\]
Every correction receives total weight one, since
$(1-D_e)+\sum_{t\ni e}p_t/3=1$. This identity is the reason to add
the same correction to both local inequalities.

Now take $H$ to be the majority graph of $k$ independently sampled
inputs. For each pair put $m=\min(w_e^+,w_e^-)$ and $u=1-2m$.
Following the sampled sign costs $m$ with probability $1-\phi_k(m)$
and $1-m$ with probability $\phi_k(m)$. Since
$\phi_k(1-m)=1-\phi_k(m)$,
\begin{align*}
 \mathbb E_H[T_k(s_e)]
 &=18u\phi_k(m)(1-\phi_k(m))
   -18u(1-\phi_k(m))\phi_k(m)=0.
\end{align*}
Taking expectation, choosing $F$ optimal, and dividing by seven proves
\eqref{eq:centered-main}.

Given $0<\tau\le1/4$, choose the smallest odd integer satisfying
\[
 k\ge87,\qquad k\ge10\log\frac{72}{7\tau}.
\]
The expected better candidate has cost at most $(11/7+\tau)\operatorname{OPT}$.
The number of clique buckets per vertex obeys
\[
 K_{\rm buck}=\binom{k}{(k+1)/2}\le2^k
 \le 2^{89}+4\left(\frac{72}{7\tau}\right)^{10\log 2}
 =\tau^{-O(1)}.
\]
The very large absolute constant arising from $k\ge87$ should be
kept in view. This result establishes polynomial accuracy dependence;
it is not intended as a practical small-constant implementation.

Using the same independent repetition and evaluation as the draft gives
$(11/7+\tau+\varepsilon)$ approximation with failure probability
$\zeta$. The bucket operations and the maintenance bounds are polynomial
in $k$ and $K_{\rm buck}$. Thus, with the hashing improvement, the offline
running time is
\[
 \tau^{-O(1)}\,n\varepsilon^{-3}
       \log^2\frac{n+2}{\varepsilon\zeta},
\]
and expected amortized implicit update time is
\[
 \tau^{-O(1)}\,\varepsilon^{-3}
       \log^2\frac{U+2}{\varepsilon\zeta}\log(N+2).
\]
Membership and same-cluster queries remain worst-case $O(1)$ under the
stated word-RAM assumptions. For $0<\varepsilon\le1/2$, choosing $\tau=\varepsilon/2$ and using
$\varepsilon/2$ for repetition and evaluation gives a
$(11/7+\varepsilon)$ approximation with polynomial dependence on
$1/\varepsilon$, near-linear offline time, and polylogarithmic dynamic
update time.

For larger requested errors, run the construction at error $1/2$.

\section{A logarithmic improvement in fully dynamic update time}
\label{sec:hashing}

Let $T$ be the number of trials, let $s$ be the evaluation-sample size,
and put $\ell_t=\log(n_t+2)$, where
$n_t=\max\{|V_{t-1}|,|V_t|\}$. In the standard word-RAM model, the
implicit fully dynamic algorithm can be implemented with expected total
time
\[
 O\!\left(Ts+T\sum_{t=1}^{u}(s\ell_t+\ell_t^2)\right)
\]
through any prefix of $u\le U$ updates. Its membership and same-cluster
queries take $O(1)$ worst-case time. Its approximation distribution and
space bound $O(Ts(n+1))$ with $n$ active vertices are preserved.

Here the word-RAM assumption is that vertex identifiers and input labels
are constant-word keys and standard constant-time universal hashing is
available. Integers up to $O(s(N+2)^2)$ and candidate-array indices
fit in a constant number of words. The bucket count is a fixed constant;
this includes every fixed-sample construction in Section~\ref{sec:centered-finite}. The random-ranking comparison convention from the draft is
retained. All dictionary randomness is independent of the samples and
pivot rankings.

The only new data-structural ingredient is dynamic perfect hashing. The
dictionary of Dietzfelbinger, Karlin, Mehlhorn, Meyer auf der Heide,
Rohnert, and Tarjan supports exact lookup in worst-case constant time,
insertion and deletion in expected amortized constant time, and linear
space. Its randomness affects running time and never changes a returned
key or value. Constant-size tuples of word keys can be encoded with a
constant-factor increase in word size. This is a Las Vegas dictionary.
See \cite{perfecthash}.

Use this dictionary for the nonzero block-intersection counters
$b_{j,p,b}$ and block sizes. Keep the balanced trees that are ordered by
owner priority and the priority queue used by pivot repair. In
particular, no replacement is made for a data structure that must support
an ordered operation. Maintain a global vertex dictionary that maps each
active vertex identifier to a stable record containing an array of its
current owner or input label in every candidate. Keep these records in a
linked list for enumeration.

Fix the full input sequence and condition on every sampled partition
index, every evaluation index, and every pivot priority. Resolve all
algorithmic ties by fixed identifiers. The complete sequence of owners,
counter keys, and dictionary requests is then independent of dictionary
randomness. Dictionary relocation changes physical locations but stable
record pointers are retained. Consequently, the expected aggregate cost
of all dictionary operations is a constant times their number, even
though that number and the sequence depend on the conditioned sampling
and pivot randomness. Averaging afterward gives the same bound without
conditioning.

For one Pivot candidate let $B_t$ count intermediate owner reports,
including insertion/deletion of the updated vertex, as in the draft.
The owner-reporting lemma gives
$\mathbb E B_t=O(\ell_t)$ and the pivot-sensitivity lemma gives
$\mathbb E F_t=O(1)$. The ordered owner-index work is
\[
 O\!\left(\sum_{t=1}^{u}(B_t+F_t+1)\ell_t\right).
\]
Its expectation is $O(\sum_t\ell_t^2)$. The pivot and certificate work
is $O(\sum_t\ell_t)$ in expectation for the constant bucket count.
Every insertion, deletion, or owner report causes $O(s)$ exact dictionary
operations and $O(s)$ arithmetic operations on score counters. Therefore
the expected aggregate score work is
\[
 O\!\left(s\sum_{t=1}^{u}(\mathbb E B_t+1)\right)
 =O\!\left(s\sum_{t=1}^{u}\ell_t\right).
\]
A sampled-input candidate costs $O(s)$ expected amortized time per
update. Summing over all trials, adding initialization, reading sampled
labels, and the $O(T)$ minimum-score selection per update proves the
claimed bound.

Let
\[
 \Lambda=\log\frac{U+2}{\varepsilon\zeta},\qquad
 T=\Theta(\varepsilon^{-1}\Lambda),\qquad
 s=\Theta(\varepsilon^{-2}\Lambda).
\]
The stream starts empty, hence $N\le U$. Since
$\varepsilon\le1$ and $\zeta<1$, we have
$s=\Omega(\log(N+2))$. Thus the ordered-index term is absorbed into the
score term and the expected amortized update time becomes
\[
 O\!\left(\varepsilon^{-3}\Lambda^2\log(N+2)\right).
\]
This saves a factor $\log(N+2)$ over the draft's update bound.

A membership query reads the selected candidate index, looks up the
vertex in the global perfect-hash dictionary, and returns the appropriate
entry of its label array. These are worst-case constant-time operations.
Two such lookups and a label comparison answer a same-cluster query.
Enumerating the linked list gives all vertex-label pairs in $O(n+1)$
time. Since dictionary operations are exact, none of these changes
introduces an approximation failure event or requires reducing $\zeta$.

\paragraph{Combined guarantees.}
The seven-sample algorithm gives simultaneous
$(1.64+\varepsilon)$ approximation with probability at least
$1-\zeta$, expected amortized update time
$O(\varepsilon^{-3}\Lambda^2\log(N+2))$, worst-case $O(1)$ membership
and same-cluster queries, and $O(n+1)$ time to list all vertex-label
pairs. The fifteen-sample variant gives $1.6+\varepsilon$ with the
same asymptotic bounds and larger constants.

The construction in Section~\ref{sec:centered-poly} gives
$11/7+\varepsilon$ approximation in
\[
 n\varepsilon^{-O(1)}\log^2\frac{n+2}{\varepsilon\zeta}
\]
offline time and
\[
 \varepsilon^{-O(1)}\log^2\frac{U+2}{\varepsilon\zeta}\log(N+2)
\]
expected amortized implicit update time. Its active space is
$(n+1)\varepsilon^{-O(1)}\log^2((U+2)/(\varepsilon\zeta))$.
The key-processing cost and bucket multiplicities are included in
the polynomial factors. Vertex-identifier lookups, stored score
arithmetic, and addressing the allocated data structures have the
word-RAM conventions stated above. Query and listing bounds are
unchanged. Section~\ref{sec:fixed-ratio-improvement} gives the
separate fixed-factor $1.64$ guarantee with no accuracy parameter.

\section{How the techniques differ from the draft}
\label{sec:technique-comparison}

Draft lemma and equation numbers refer to the compiled \texttt{main.tex}.
The explicit dynamic lemmas appear in Appendix~B.

\paragraph{Weighted online history bound.}
Draft Lemma~3.1 (History bound for a triple candidate) bounds the
partition distance in Eq.~(13) and then charges $\eta+1$ for every pair
on which the majority graph disagrees with the reference partition $F$.
Lemma~\ref{lem:improved-weighted-history} bounds the consensus cost
difference directly. If the candidate reverses $F$ on pair $e$, its cost
changes by $1-2p_e$, where $p_e$ is the fraction of inputs disagreeing
with $F$ on that pair. Adding the solver's error bound on the signed
graph gives a charge of $\eta+1-2p_e$ per majority error. The bound
applies at every earlier prefix against the restriction of $F$, so it
also controls the maximum over those prefixes.

\paragraph{Improved online competitive ratio.}
Draft Lemma~3.2 (Expected history score) uses three distinct sampled
inputs and the bound $3p_e^2$ on a majority error. Together with draft
Lemma~3.1, this gives a majority-candidate history bound of
$\operatorname{OPT}_t+3(\eta+1)X_t$, where $X_t=\sum_e p_e^2$ for an
optimal reference partition. Lemma~\ref{lem:improved-majority} uses
five or fifteen samples drawn with replacement and bounds the \emph{weighted}
error by $2\eta p_e^2$. The majority-candidate bound therefore becomes
$\operatorname{OPT}_t+2\eta X_t$. The input candidate still has expected
cost $2\operatorname{OPT}_t-2X_t$, as in draft Eq.~(14).
Theorem~\ref{thm:improved-online-exact} balances the two bounds with input
probability $\eta/(\eta+1)$ and uses the amplification and phase selection
of draft Lemmas~3.3--3.4. This changes the ratio from
$2-2/(3\eta+5)+\varepsilon$ to $2-1/(\eta+1)+\varepsilon$ with the same
Online Agreement solver. When $\eta\ge16/3$, the variant in
Section~\ref{sec:improved-online} obtains the same weighted error bound
using only majorities of five or seven samples.

\paragraph{Online recourse independent of the input count.}
Draft Lemma~3.5 (Per-vertex recourse) starts the phase count at $1/m$,
the smallest possible positive exact score.
Theorem~\ref{thm:improved-online-estimated} uses the minimum of independent
group means from draft Lemmas~B.6--B.8 with the deterministic Online
Agreement candidates. Draft Lemma~B.6 compares a candidate and $F_t^*$
at the same prefix $t$. Equation~\eqref{eq:improved-eval-comparison}
also compares a candidate at $u$ with $F_t^*|_{V_u}$ for every $u\le t$.
These comparisons let the weighted bound from
Lemma~\ref{lem:improved-weighted-history} control the estimated history.
The grid of scores has spacing $1/g$, where $g$ is the group size, so
the logarithm in the phase count depends on $g$ in place of $m$.

\paragraph{Fewer switches within an online phase.}
Draft Lemma~3.5 allows up to $K$ outer switches per phase because
eligibility is monotone within a phase.
Theorem~\ref{thm:improved-random-recourse} uses a fresh random order in
each phase and chooses the first eligible candidate in that order.
Given the elimination order, the number of selections is bounded by
the number of record minima. Its mean is at most
$H_K=O(\log(K+1))$ per phase, and its exponential moment gives a tail
bound for the total across phases. The selection inequality from draft
Lemma~3.4 still applies to any eligible choice. Changes within the
solver copies are counted separately as in draft Lemma~3.5.

\paragraph{Conditional transfer to explicit dynamic consensus.}
Draft Lemma~B.1 uses the triple candidates in the reduction, and
Lemmas~B.5--B.8 combine the proxy, solver success, and estimated score
comparison. Section~\ref{sec:conditional} substitutes the weighted
proxy based on Lemmas~\ref{lem:improved-majority}
and~\ref{lem:improved-weighted-history}, so the reduction gives
$2-1/(\eta+1)+\varepsilon$ under the solver hypothesis. In the draft,
the pairwise intersections of the three buckets coincide, which allows
sampling a uniform neighbor in $O(\log(n+2))$ time. The larger majorities
use inclusion-exclusion over ordered buckets to count neighbors in an
arrival-index prefix. Binary search for a uniform rank takes
$O(\log^2(n+2))$ time, which fits the polylogarithmic factor in the
update bound of draft Lemma~B.1.

\paragraph{A better bound for the original seven-sample algorithm.}
Draft Lemma~4.3 (Exact pair identity) expresses the probability $D_e$
that Pivot opposes the sampled sign as a sum of bad-triangle weights.
Draft Lemmas~4.4--4.5 then bound positive corrections caused by sampling,
including the saving for unanimous errors. Draft Lemmas~4.6--4.7 use
the resulting sampling penalty in their pair and triangle bounds for
the full majority graph. Equations~\eqref{eq:finite-centered-pair}
and~\eqref{eq:finite-centered-triplet} instead use the sampled sign costs
and add the same term $\lambda t_k(s_e)$ to both bounds. Draft
Lemma~4.3 gives that term total weight one, and
Lemma~\ref{lem:centered-zero} cancels it in expectation. For $k=7$,
the new certificates give expected factor $41/25$ in
Theorem~\ref{thm:centered-finite}. The repetition and evaluation steps
from draft Theorem~4.1 then give $1.64+\varepsilon$ in place of
$1.6944+\varepsilon$ with the original 35 buckets.

\paragraph{Approximation tradeoffs at other fixed sample sizes.}
Draft Lemma~4.6 certifies a pair inequality and a triangle inequality
for seven samples. Theorem~\ref{thm:centered-finite} and its finite-sample
extension certify \eqref{eq:finite-centered-pair}--\eqref{eq:finite-centered-triplet}
with different $(\beta,\lambda)$ at each sample size. Five samples
give $5/3+\varepsilon$ with ten buckets. Nine, thirteen, and fifteen
samples give $1.625+\varepsilon$, $1.61+\varepsilon$, and
$1.6+\varepsilon$ with 126, 1716, and 6435 buckets, respectively.
The sampled-majority graph uses $\binom{k}{(k+1)/2}$ buckets per vertex
in the representation analyzed by draft Lemma~4.2, so these choices
change the constants in the time and space bounds.

\paragraph{A fixed factor of $1.64$.}
The proof of draft Theorem~4.1 amplifies the single-trial bound in
Eq.~(22) and estimates candidate costs with accuracy set by
$\varepsilon$. Section~\ref{sec:fixed-ratio-improvement} applies those
steps to the nine-sample expected factor $13/8$. With fixed slack
$\delta=1/200$, amplification and evaluation give
$(1+\delta)(13/8+\delta)=1.63815<1.64$. The gap between $13/8$ and $1.64$
thus lets us fix the accuracy used to set the trial count and
evaluation sample size.

\paragraph{The $11/7+\varepsilon$ guarantee.}
Draft Lemmas~4.6--4.7 prove a bound for a fixed sample size. For growing
odd $k$, \eqref{eq:centered-pair}--\eqref{eq:centered-triplet} replace
those local inequalities using $18t_k(s)$ in the charges scaled by
seven. The potentially tight triangle has sign costs $(1/4,1/4,1/2)$.
Equation~\eqref{eq:centered-critical-bound} absorbs the corrections on
the first two edges into quadratic slack up to $O(e^{-k/10})$, and the
third edge uses $18t_k(r)\le9-18r$. Summing with the weights from draft
Lemma~4.3 and applying Lemma~\ref{lem:centered-zero} gives the
single-trial bound \eqref{eq:centered-main}. The loss above $11/7$ is
$O(e^{-k/10})$, improving the $O(k^{-1/2})$ loss in earlier versions
of these notes. Theorem~\ref{thm:eleven-sevenths} reaches
$11/7+\varepsilon$ with $k=O(\log(2/\varepsilon))$ and a bucket count
polynomial in $1/\varepsilon$.

\paragraph{Faster implicit score maintenance.}
Draft Lemmas~4.8--4.12 supply exact Pivot maintenance, its time bound,
and the expected number of owner reports. Draft Lemma~4.13 (Maintaining
a candidate and its score) performs $O(s)$ counter operations per
report, paying $O(\log(n_t+2))$ for each one in a balanced dictionary.
Section~\ref{sec:hashing} uses exact dynamic perfect hashing for the
counters and retains the ordered owner indexes. With
$\ell_t=\log(n_t+2)$ as defined there, the expected amortized score
work per update drops from $O(s\ell_t^2)$ to $O(s\ell_t)$ for a fixed
bucket count. The separate $O(\ell_t^2)$ ordered-index work fits within
$O(s\ell_t)$ because $s=\Omega(\log(N+2))$. This removes one factor
of $\log(N+2)$ from the update bound in draft Theorem~2.2.

\paragraph{Constant-time implicit queries.}
The membership lookup in the proof of draft Theorem~2.2 searches the
selected candidate's balanced dictionary in $O(\log(n+2))$ time for
$n$ active vertices. Section~\ref{sec:hashing} uses one perfect-hash
lookup to find the vertex record and an array access at the selected
candidate index to read its label. This gives worst-case $O(1)$
membership on a word RAM. Two lookups and one label comparison give
the same bound for same-cluster queries.

\begin{thebibliography}{9}
\bibitem{acn}
Nir Ailon, Moses Charikar, and Alantha Newman.
\emph{Aggregating Inconsistent Information: Ranking and Clustering}.
Journal of the ACM 55(5), Article 23, 2008.
\url{https://doi.org/10.1145/1411509.1411513}.

\bibitem{online}
Vincent Cohen-Addad, Silvio Lattanzi, Andreas Maggiori, and Nikos Parotsidis.
\emph{Online and Consistent Correlation Clustering}.
Proceedings of ICML, PMLR 162, 4157--4179, 2022.
Theorems E.2 and F.1.
\url{https://proceedings.mlr.press/v162/cohen-addad22a.html}.

\bibitem{perfecthash}
Martin Dietzfelbinger, Anna Karlin, Kurt Mehlhorn, Friedhelm Meyer auf
der Heide, Hans Rohnert, and Robert E. Tarjan.
\emph{Dynamic Perfect Hashing: Upper and Lower Bounds}.
SIAM Journal on Computing 23(4), 738--761, 1994.
\url{https://doi.org/10.1137/S0097539791194094}.
\end{thebibliography}
\end{document}
